% GPT-6 Astra Pro; effort not exposed; /root/replacement; 22 November 2026.
\section[HD-F1-R19: rational rank-one action rigidity]{HD-F1-R19: rational rank-one action rigidity and adjacent objects with a one-dimensional invariant cross-Hom}
\label{sec:followup-rm-rank-one-action}

Use the genuine real-quadratic fourfold, rational algebraic
$B$, compatible polarization, and balanced character of
R07--R11. The following actual adjacent class contains no
weakly semiregular object at virtual rank $1$ or $-1$:
\[
 \begin{gathered}
 H^0(E)=A,\quad H^1(E)=T,\quad H^j(E)=0\ (j\ne0,1),\\
 \operatorname{ch}(E)=N e^{-B}v_q,\quad
 v_q=(1-qA_1)(1-qA_2),\quad q>0,\quad N\ne0,\\
 \operatorname{Ext}^{j}_G(E,E)=0\ (j<0),\qquad
 h:=\dim\operatorname{Hom}_G(A,T)\le1,
 \end{gathered}
 \tag{F1R19.1}
\]
where $A,T$ are positive-rank torsion-free sheaves,
slope semistable for one ample polarization, and $G$
is any specified coherent faithful finite point action
with an actual lift. Every actual connecting class is
included. This extends R18 to a nonzero one-dimensional
invariant cross-Hom group. We prove the exact negative-Ext
map, the full endpoint characters, and a stronger integral
inequality for arbitrary positive-rank coherent endpoints.

\paragraph{1. The exact exterior model on the complete HKR space.}
Work first after the explicit rational $B$-shear of R07.
Write $U=U_1\oplus U_2$, where
$U_i=T_i\oplus W_i$ has dimension four. On the $i$th
block take generators $x_1,x_2,y_1,y_2$ as in R11,
with $A_i=x_1x_2y_1y_2$, and let
$t_a=\iota_{x_a^*}$ and $w_b=y_b\wedge$. The map
\[
 \Phi_i:\bigwedge^*U_i\longrightarrow
    \bigwedge^*\langle x_1,x_2,y_1,y_2\rangle,
       \qquad a\longmapsto D_a(x_1x_2)
 \tag{F1R19.2}
\]
is an isomorphism of modules over $\bigwedge^*U_i$.
Indeed its sixteen exterior basis monomials remove
an arbitrary subset of the two $x$'s and add an
arbitrary subset of the two $y$'s; the resulting
sixteen form monomials are distinct and have nonzero
coefficients. Anticommutation proves the module
identity. It identifies exterior degree $j$ with
Hochschild homological degree $2-j$.

Tensoring the two actual maps in (F1R19.2), with the
exterior signs, therefore gives
\[
       HH_0(X)\simeq\bigwedge^4U,\qquad
       HH_{-2}(X)\simeq\bigwedge^6U,
 \tag{F1R19.3}
\]
under which the action of $K=K_1\oplus K_2$ is
exterior multiplication. In the displayed convention
\[
 \xi_i=-t_1t_2-qw_1w_2\in\bigwedge^2U_i,
   \qquad \Phi_i(\xi_i)=1-qA_i,\qquad \xi_i^2\ne0.
 \tag{F1R19.4}
\]
Thus $K_i=\{a\in\bigwedge^2U_i:a\xi_i=0\}$.
The wedge pairing on $\bigwedge^2U_i$ is nondegenerate,
$\xi_i$ is nonisotropic for it, and
\[
                   \bigwedge^2U_i=\mathbf C\xi_i\oplus K_i
 \tag{F1R19.5}
\]
is an orthogonal decomposition with nondegenerate
restriction to $K_i$.

\paragraph{2. Rank at most one over the complex numbers.}
For $w\in\bigwedge^4U$ decompose
\[
 w=\sum_{j=0}^4 w_j,\qquad
 w_j\in\bigwedge^jU_1\otimes\bigwedge^{4-j}U_2.
 \tag{F1R19.6}
\]
We classify the vectors for which the map
$K\to\bigwedge^6U$, $a\mapsto aw$, has rank at
most one. A nonzero $w_0$ makes the projection of
the $K_1$ action to bidegree $(2,4)$ have rank five.
A nonzero $w_4$ gives the same conclusion for $K_2$
and bidegree $(4,2)$. Consequently both vanish.

A nonzero $w_1$ has a nonzero coefficient
$a\in U_1$ after applying a linear functional to
its $\bigwedge^3U_2$ factor. The projected $K_1$
action then contains the map
\[
                    K_1\longrightarrow\bigwedge^3U_1,
                               \qquad k\longmapsto ka.
 \tag{F1R19.7}
\]
On all of $\bigwedge^2U_1$ this map has rank three
and kernel $a\wedge U_1$, of dimension three.
The functional $k\mapsto k\xi_1$ is nonzero on
$a\wedge U_1$: the pairing
$(a,b)\mapsto a b\xi_1$ on $U_1$ is nondegenerate
because $\xi_1^2\ne0$. Hence the intersection of
that kernel with $K_1$ has dimension two, and
(F1R19.7) still has rank three. This excludes $w_1$.
The identical argument with the two actual factors
interchanged excludes $w_3$.

It remains to consider
$w=w_2\in\bigwedge^2U_1\otimes\bigwedge^2U_2$.
By (F1R19.5) write it uniquely as
\[
 w=a\xi_1\xi_2+b_1\xi_2+\xi_1b_2+M,
       \qquad b_i\in K_i,\quad M\in K_1\otimes K_2.
 \tag{F1R19.8}
\]
The $K_1$ and $K_2$ action images lie in the distinct
bidegrees $(4,2)$ and $(2,4)$, so their ranks add.
At total rank at most one at least one action is zero.
If the $K_2$ action is zero, its nondegenerate wedge
pairing and (F1R19.8) give $b_2=M=0$. If the $K_1$
action is zero, they give $b_1=M=0$. Thus the complete
complex rank-at-most-one locus is
\[
 \left(\bigwedge^2U_1\right)\otimes\mathbf C\xi_2
 \quad\cup\quad
 \mathbf C\xi_1\otimes\left(\bigwedge^2U_2\right).
 \tag{F1R19.9}
\]
Conversely a vector in either subspace has just one
possibly nonzero action, landing in a one-dimensional
image, so it does have rank at most one. This proves
both inclusions, on all seventy dimensions in (F1R19.3).

\paragraph{3. Rational characters force the balanced line.}
Let $z$ be a rational class in
$\bigoplus_p H^{p,p}(X)$ whose $K$-action has rank
at most one. The balanced vector $v_q$ is rational:
$e^B\operatorname{ch}(E)=Nv_q$ and $B$ is rational.
Under the first alternative in (F1R19.9) one has
\[
       z=(r+\ell_1+dA_1)(1-qA_2),
 \tag{F1R19.10}
\]
where $r$ is its rational rank, $\ell_1$ is of type
$(1,1)$ entirely in the first RM block, and $d$ is
a complex scalar. In particular the codimension-one
component of $z$ is $\ell_1$, and is rational.

Choose a positive integer $m$ clearing the denominator
of $1+\sqrt{\Delta_0}$ in the actual real multiplication
algebra, so
\[
               j=m(1+\sqrt{\Delta_0})\in\operatorname{End}(X).
 \tag{F1R19.11}
\]
Its pullback is a rational linear map on every rational
cohomology group. On the first Hodge block it has
eigenvalue $\alpha_1=m(1+\sqrt{\Delta_0})$.
Therefore $j^*\ell_1=\alpha_1^2\ell_1$, where
$\alpha_1^2$ is irrational. A nonzero rational vector
cannot satisfy this equation: in any rational basis,
a nonzero rational coordinate would express
$\alpha_1^2$ as the ratio of two rational coordinates.
It follows that $\ell_1=0$.

Now the codimension-two part of $z-rv_q$ is
$(d+rq)A_1$ and is rational. Its $j^*$ eigenvalue is
$\alpha_1^4$, also irrational, since its
$\sqrt{\Delta_0}$ coefficient is
$4m^4(1+\Delta_0)\ne0$. The same rational-coordinate
argument gives $d+rq=0$. Hence $z=rv_q$.
For the second alternative in (F1R19.9), the identical
proof uses the conjugate eigenvalue
$\alpha_2=m(1-\sqrt{\Delta_0})$; its square and fourth
power are again irrational. We have proved the exact
rational rigidity assertion
\[
 z\in\bigoplus_p H^{p,p}(X)\cap H^{\mathrm{even}}(X,\mathbf Q),
 \quad\operatorname{rank}(K\xrightarrow{a\mapsto az}HH_{-2})\le1
       \quad\Longrightarrow\quad z=\operatorname{rk}(z)v_q.
 \tag{F1R19.12}
\]
The formula includes rank-zero classes, where its
conclusion is $z=0$.

To return to the given character, set
$K_E=\ker c^2_{\operatorname{ch}(E)}$.
The actual shear $s_B$ of (F1R7.5) carries $K_E$ to
$K$ and identifies $c_z^2|_{K_E}$ with
$e^{-B}c_{e^Bz}^2|_K$. All these maps are invertible
on their source or target, and $e^Bz$ is rational.
Thus (F1R19.12) applies without changing the rank of
the given restriction. This explains the exact
normalization morphism; no change of object is presumed.

\paragraph{4. Gluability and the actual endpoint Yoneda kernels.}
The lifted canonical triangle is
\[
           A\xrightarrow{i}E\xrightarrow{p}C=T[-1]
                                  \xrightarrow{\delta}A[1].
 \tag{F1R19.13}
\]
It first gives the exact isomorphism
\[
 \operatorname{Hom}_G(T,A)\xrightarrow{\sim}
       \operatorname{Ext}^{-1}_G(E,E),\qquad
       f\longmapsto i[-1]f[-1]p,
 \qquad \operatorname{Ext}^{j}_G(E,E)=0\ (j<-1).
 \tag{F1R19.14}
\]
Indeed $\operatorname{Hom}(E,C[-1])$ and
$\operatorname{Hom}(E,C[-2])$ vanish by the standard
t-structure, so the first shifted Hom sequence
identifies $\operatorname{Hom}(E,A[-1])$ with
$\operatorname{Hom}(E,E[-1])$. The other sequence
identifies the former with
$\operatorname{Hom}(C,A[-1])=\operatorname{Hom}(T,A)$:
its adjacent terms are $\operatorname{Ext}^{-2}(A,A)$
and $\operatorname{Ext}^{-1}(A,A)$, both zero.
For $j<-1$ the same sequences reduce to
$\operatorname{Ext}^{j+1}(T,A)=0$. Taking the invariant
subspaces preserves these exact statements.
Consequently the gluability in (F1R19.1) is exactly
$\operatorname{Hom}_G(T,A)=0$.

Weak semiregularity, with the actual identity-sector
comparison of R05, gives
$\operatorname{ev}_E^2(K_E)=0$. The exact two Hom
sequences (F1R18.4), now retaining their preceding
terms, imply
\[
 \begin{split}
 \operatorname{ev}_A^2(K_E)&\subset
       \delta\circ\operatorname{Hom}_G(A,T),\\
 \operatorname{ev}_T^2(K_E)&\subset
       \operatorname{Hom}_G(A,T)[2]\circ\delta.
 \end{split}
 \tag{F1R19.15}
\]
Here $\delta$ is equivalently the actual class
$T\to A[2]$. Naturality along $i,p$ gives precisely
these two kernel memberships; no vanishing of the
ordinary cross-Hom groups is required. Taking the
ordinary full trace shows that both
$c^2_{\operatorname{ch}(A)}|_{K_E}$ and
$c^2_{\operatorname{ch}(T)}|_{K_E}$ have rank at most
$h\le1$. Their rational characters therefore satisfy
\[
          \operatorname{ch}(A)=r_Ae^{-B}v_q,
      \qquad \operatorname{ch}(T)=r_Te^{-B}v_q,
      \qquad N=r_A-r_T.
 \tag{F1R19.16}
\]
This conclusion does not assert that the endpoint
characteristic Ext classes in (F1R19.15) themselves
vanish when $h=1$; only their traced action and their
full characters have been determined.

\paragraph{5. The corrected Euler inequality with nonzero forward Hom.}
Retain the actual notation and maps from R18:
\[
 \begin{gathered}
 d=\dim\operatorname{End}_G(A)+\dim\operatorname{End}_G(T),\quad
 V=\operatorname{Ext}^2_G(T,A),\\
 a=\operatorname{ext}^1_G(T,A),\quad
 b=\operatorname{ext}^1_G(A,T),\quad c=\chi_G(T,A),\\
 d_\delta(u,v)=u[2]\delta-\delta v,\quad
 s=\dim\ker d_\delta,\quad L_\delta=\operatorname{im}d_\delta.
 \end{gathered}
 \tag{F1R19.17}
\]
The actual endomorphism sequence (F1R18.8) and the
exact tangent kernel (F1R18.10) use only negative Hom
vanishings of sheaves, and still give
\[
                  e_0=a+s,\qquad
       V/L_\delta\lhook\joinrel\longrightarrow
          \operatorname{Ext}^1_G(E,E),\quad
                         v\longmapsto i[1]vp.
 \tag{F1R19.18}
\]
For completeness, the first precomposition
$V\to\operatorname{Hom}_G(E,A[1])$ has kernel
$\operatorname{End}_G(A)[2]\delta$; the next
postcomposition has kernel induced by
$\operatorname{Hom}_G(E,C)=\operatorname{End}_G(T)p$.
Their combined inverse image is exactly $L_\delta$.
The compatible endpoint pairs in the endomorphism
sequence are exactly $\ker d_\delta$ and its first
kernel is $\operatorname{Ext}^1_G(T,A)$, including
$\delta=0$.

The eight point-action tangent directions are disjoint
from that quotient image. Composition with $i$ kills
$i[1]vp$, while the natural endpoint tangent has the
injective postcomposition map (F1R18.11), whose
preceding term is the negative group
$\operatorname{Hom}(A,C)=0$. Equation (F1R19.16)
makes the endpoint degree-one character action rank
eight. Thus the same actual composition argument gives
\[
                         e_1\ge8+\dim V-d+s.
 \tag{F1R19.19}
\]
Gluability and invariant CY4 duality give
$\operatorname{Ext}^4_G(T,A)\simeq
 \operatorname{Hom}_G(A,T)^*$, of dimension $h$,
and $\operatorname{Hom}_G(T,A)=0$. Hence the exact
cross Euler identity is now
\[
               \dim V=c+a+b-h,\qquad
               e_2=18,\qquad k:=\chi_G(E,E)=2e_0-2e_1+18.
 \tag{F1R19.20}
\]
Combining (F1R19.18)--(F1R19.20) cancels $a$ and $s$
and proves
\[
                         c+b+\frac{k}{2}\le d+1+h.
 \tag{F1R19.21}
\]
The exact induction/forgetful Euler pairing, polarized
on $A,T,A\oplus T$ as in R18, and the actual endpoint
characters (F1R19.16) give
\[
 k=\frac{4Dq^2N^2}{|G|}\in2\mathbf Z_{>0},\qquad
 c=\frac{k r_A r_T}{N^2}\in\mathbf Z_{>0},\qquad
              \frac{k(r_A^2+r_T^2)}{2N^2}+b\le d+1+h.
 \tag{F1R19.22}
\]
No cross-Hom vanishing has been inserted in that Euler
pairing; its new degree-four term is exactly the
$h$ retained in (F1R19.20).

In the reduced positive coprime unequal ranks
$r_A=gp$, $r_T=gq_0$, the same divisibility as R18
gives a positive integer $t$ with
\[
 k=t(p-q_0)^2,\qquad c=tpq_0,\qquad
 d+1+h\ge\frac{t(p^2+q_0^2)}2+b\ge5+b.
 \tag{F1R19.23}
\]
The last inequality follows from the parity of
$t(p-q_0)^2$: for odd $p-q_0$ one has $t\ge2$ and
$p^2+q_0^2\ge5$, whereas for even $p-q_0$ the two
coprime ranks are odd and their squared sum is at
least ten. Thus the actual coherent endpoint class
in (F1R19.1) always satisfies
\[
                                  d+h\ge4+b.
 \tag{F1R19.24}
\]
In particular invariant-simple endpoints, with $d=2$,
are excluded even at $h=1$, in every nonzero virtual
rank. For torsion-free endpoints of virtual rank
$|N|=1$, S-22 and (F1R19.16) give both ranks at least
two, so the reduced squared sum is at least thirteen
and $t\ge2$. The stronger exact restriction is then
\[
                                  d+h\ge12+b.
 \tag{F1R19.25}
\]
This excludes every such pair with $d+h\le11$ without
any semistability assumption.

Finally, for the slope-semistable torsion-free endpoints,
the actual filtration and reflexive-Hom proof in R18
applies to their now established full characters. Put
$R=r_A^2+r_T^2$ and $S=r_A+r_T$. It gives
\[
                   d+1+h\le R-2S+5+h<R,
 \tag{F1R19.26}
\]
since $r_A,r_T\ge2$ and $h\le1$. Equation
(F1R19.22) therefore implies
\[
                              2\le k\le2N^2-2.
 \tag{F1R19.27}
\]
This proves the stated exclusion for $|N|=1$ and also
excludes the minimal profile $(e_0,e_1,e_2)=(3,8,18)$,
whose Euler number is eight, at $|N|\le2$ in this
adjacent class. In particular every gluable weakly
semiregular virtual-rank-one adjacent object with
semistable torsion-free endpoints must have
$\dim\operatorname{Hom}_G(A,T)\ge2$.

\paragraph{Framework, sources, and remaining scope.}
The map (F1R19.2) is the explicit cyclic exterior-module
map for the HKR action. The characteristic/trace
identity is precisely Buchweitz--Flenner's published
Proposition~6.2.1, Corollary~6.3.2 and Theorem~6.4.2,
with the stated $B$-shear; these complete relevant proofs
were checked in the earlier R07--R11 source events.
The rationality step uses an actual denominator-cleared
RM endomorphism on rational cohomology. The endpoint
maps are the exact canonical truncation triangle and
its two Hom sequences, and the integral Euler pairing
uses the same actual induction/forgetful maps as S-21
and the published Hochschild Riemann--Roch comparison
checked for R16--R18. The endomorphism estimate is the
proved coherent-filtration bound (F1R18.25), not an
assumption that endomorphisms preserve that filtration.

Fresh searches on Hexagon, arXiv and journal sites were
made for the rank-one exterior-action and equivariant
semiregularity questions. They provide comparison only;
no unrefereed theorem is used in this proof and no
exhaustive originality claim is made. This argument
determines the rank-at-most-one rational locus and the
specified adjacent class. It constructs no higher-rank
seed and makes no assertion excluding every adjacent
object with larger invariant forward Hom, nonsemistable
endpoints, or arbitrary cohomological amplitude.
